\documentclass{book}
\usepackage{template}
 
 \begin{document}

% Aparentemente o Tikz tem um problema de escurecimento das cores RGB, então temos que definí-las primeiro
\definecolor{Primeira}{RGB}{255, 153, 254}
\definecolor{Segunda}{RGB}{255, 205, 205}
\definecolor{Terceira}{RGB}{255, 205, 254}
\definecolor{Quarta}{RGB}{129, 204, 0}
\definecolor{Quinta}{RGB}{0, 204, 254}
\definecolor{Sexta}{RGB}{254, 204, 151}
% Dentro dos desenhos tentei usar o comando de fração, mas como o Tikz não entende dfrac, apenas o frac
% ficou muito pequeno. Portanto, utilizei o "draw" para fazer os números
 
\begin{tikzpicture}[x=1mm, y=1mm]
% Fita principal (metade cinza, metade branca)
\draw[thick] (0,0) rectangle (100,10);
\draw[thick, fill=lightgray] (0,0) rectangle (100/2,10);
\end{tikzpicture}
 
\begin{tikzpicture}[x=1mm, y=1mm]
% Fita rosa de largura 1/2
\draw[thick, fill=Primeira] (0,0) rectangle (100/2,10);
\draw (100/4,7) node [] {1};
\draw (100/4,3) node [] {2};
\draw (100/4-2,5.3) -- (100/4+2,5.3);
\end{tikzpicture}
 
\begin{tikzpicture}[x=1mm, y=1mm]
% Fita rosa de largura 1/4
\draw[thick, fill=Segunda] (0,0) rectangle (100/4,10);
\draw (100/8,7) node [] {1};
\draw (100/8,3) node [] {4};
\draw (100/8-2,5.3) -- (100/8+2,5.3);
\end{tikzpicture}
 
\begin{tikzpicture}[x=1mm, y=1mm]
% Fita rosa de largura 1/6
\draw[thick, fill=Terceira] (0,0) rectangle (100/6,10);
\draw (100/12,7) node [] {1};
\draw (100/12,3) node [] {6};
\draw (100/12-2,5.3) -- (100/12+2,5.3);
\end{tikzpicture}
 
\begin{tikzpicture}[x=1mm, y=1mm]
% Fita rosa de largura 1/8
\draw[thick, fill=Quarta] (0,0) rectangle (100/8,10);
\draw (100/16,7) node [] {1};
\draw (100/16,3) node [] {8};
\draw (100/16-2,5.3) -- (100/16+2,5.3);
\end{tikzpicture}
 
\begin{tikzpicture}[x=1mm, y=1mm]
% Fita rosa de largura 1/10
\draw[thick, fill=Quinta] (0,0) rectangle (100/10,10);
\draw (100/20,7) node [] {1};
\draw (100/20,3) node [] {10};
\draw (100/20-2,5.3) -- (100/20+2,5.3);
\end{tikzpicture}
 
\begin{tikzpicture}[x=1mm, y=1mm]
% Fita rosa de largura 1/16
\draw[thick, fill=Sexta] (0,0) rectangle (100/16,10);
\draw (100/32,7) node [] {1};
\draw (100/32,3) node [] {16};
\draw (100/32-2,5.3) -- (100/32+2,5.3);
\end{tikzpicture}
  
 
 \end{document}